\begin{rSection}{Scientific Collaborations}

\begin{rSubsection}{Rubin Observatory \ac{lsst}}{2023 - Present}{}{}
I hold roles in both the \href{https://rubinobservatory.org}{Rubin Project} and the \href{https://lsstdesc.slac.stanford.edu}{\ac{lsst} \ac{desc}}.

\textit{Rubin Observatory Project Team}: Led \ac{LSSTCam} focal plane commissioning to Rubin Observatory on Cerro-Pachón in Chile, from receipt on summit through early science operations. Led testing of the focal plane in-laboratory and on-telescope. Characterized defects in the \ac{LSSTCam} focal plane CCDs, and developed custom image processing tool to monitor the stability of defects in the \ac{LSSTCam} focal plane. Led the \ac{ToO} observing program during Rubin early operations and first year of the \ac{lsst}.

\textit{\ac{lsst} \ac{desc}}: Forecasted the constraining power of dark sirens using simulated \ac{lsst} data under different conditions. Leading a dark-siren analysis using \ac{lsst} commissioning data (Data-Preview 2) and public data from the International Gravitational Wave Network (IGWN), the first dark siren measurement using \ac{lsst} data.
\end{rSubsection}

\begin{rSubsection}{\ac{igwn}}{2025 - Present}{}{}

Led dark siren analysis using \ac{DELVE} survey data for \ac{igwn} O4c data, including galaxy catalog preparation and systematic-uncertainty quantification. Coordinated time-critical \ac{ToO} observations with \ac{igwn} researchers, interfacing with Rubin Observatory alert brokers to rapidly disseminate information about \ac{GW} candidates. Contributed to dark siren analysis using Dark Energy Survey (DES) data for \ac{igwn} O4b data.

\end{rSubsection}

\begin{rSubsection}{\href{https://www.darkenergysurvey.org/}{\ac{des}}}{2023 - Present}{}{}
Performed observations, image processing, and analysis of images from \ac{DECam} observations in response to \ac{GW} triggers from the \ac{igwn} \ac{O4} for the \ac{desgw} observing program. Implemented improvements to \ac{desgw} alert and image processing pipelines, including improvements to the telescope strategy and \ac{GW} alert processing. Led collaboration between the \ac{J-GEM} and \ac{desgw} through joint observations of \ac{GW} event S250328ae using \ac{DECam} and the \ac{PFS} at Subaru Observatory.

\end{rSubsection}


\begin{rSubsection}{Robotic Positioners for Future Spectroscopic Surveys}{2023 - Present}{}{}
Tested Dark Energy Spectroscopic Instrument (DESI) robotic positioners to increase movement precision and improve low performing robotic positioners on the DESI focal plane. Characterized a common failure mode of DESI positioners and resolved the failure with a reproducible solution. Authored and submitted an internal report and presented the results of this study to the DESI focal plane working group. Led dynamic testing of prototype robotic positioner performance through lifetime, focal-ratio degradation, and fiber angle tests using different telescope configurations. Tested prototype robotic positioners for Stage-V surveys, focusing on characterization of tilt induced by different orientations.
\end{rSubsection}

\begin{rSubsection}{\href{https://ait.mit.edu/instruments/llamas/}{\ac{LLAMAS}}}{2020 - 2022}{}{}
Assembled the opto-mechanical mounts for the \ac{LLAMAS} instrument spectrograph. Led verification and testing of mechanical tolerance for custom-fabricated parts with a coordinate measurement machine. Verified spectrograph grating efficiency through optical tests, and integrated the spectrograph camera lenses into their bezels for the 24 camera spectrograph. Integrated the 2,400 AR-coated fibers to the spectrograph by bonding with optical adhesive, with zero fiber loss. LLAMAS was integrated to the Magellan Baade telescope, and achieved first light in November 2024.
\end{rSubsection}


\end{rSection}
